\section{A Gaussian test that cancels the affine volume}
\label{sec:gaussian}

We prove Theorem~\ref{thm:gaussian-block}. Fix its dimensions
$k=m=2\lfloor d/16\rfloor$, $t=m+1$, and $\ell=4k$, and let
$p=f\sigma_d$ with $0\le f\le L$. Draw the block data $(A,Y)$, where
$Y=AS$, and then draw $t$ independent posterior replicas conditional on
those data. Write $Q$ for the law of the replicas together with $A$,
and use $D,F,J$ for their difference matrix, difference span, and affine
volume as in \eqref{eq:equal-label-geometry}.

For large $d$, $k+m\le d-2$. Corollary~\ref{cor:replica-likelihood} applies
with
\[
 c=(2\pi)^{-km/2},\qquad
 R(d\mathbf s,dA)=p^{\otimes t}(d\mathbf s)\gamma_{k,d}^{F}(dA).
\]
The evaluated density is the common exact version of
Theorem~\ref{thm:equal-label-measure}. If $K$ is the total correlation
of the replicas, Lemma~\ref{lem:replica-logs} gives
\begin{equation}\label{eq:gaussian-correlation-upper}
 K\le\log c-k\E_Q\log J-m\E_Q\log p_A(Y)<\infty.
\end{equation}
The potentially large term is the negative logarithm of $J$. We recover
it from a test determined by the projected replicas. After that cancellation,
only a comparison between two equal-label projection densities remains.

\subsection{The observable projected test}

Let $B$ be the independent $\ell$-row matrix in
Theorem~\ref{thm:gaussian-block}, and set $U_i=BS_i$. The projected tuple
determines
\[
 BD=[U_2-U_1,\ldots,U_t-U_1].
\]
At its almost sure full rank, choose a $k\times\ell$ matrix $C_0$ with
orthonormal rows in the orthogonal complement of the columns of $BD$,
measurably as a function of $BD$. Projecting the standard basis and taking
the first successful Gram--Schmidt choices supplies such a rule.
There is room for these rows because $\ell-m=3k$. Define the selector
arbitrarily on rank failures. Put
\[
 \widetilde A=C_0B.
\]
This matrix is known from $(B,U_1,\ldots,U_t)$ and satisfies
$\widetilde A S_i=C_0U_1$ for every replica. It supplies an observable
equal-label likelihood to compare with the actual likelihood in
\eqref{eq:gaussian-correlation-upper}.

\begin{lemma}[Projected likelihood test]\label{lem:gaussian-test}
Let $G$ be an $\ell\times m$ standard Gaussian matrix and set
\[
 Z_{k,m}=\E\det(G^{\mathsf T}G)^{-k/2}.
\]
This constant is finite and positive. The nonnegative function
\begin{equation}\label{eq:gaussian-projected-test}
 \Lambda(B,U_1,\ldots,U_t)
 =\frac{c}{Z_{k,m}}
  \det((BD)^{\mathsf T}(BD))^{-k/2}
  p_{C_0B}(C_0U_1)^{-m},
\end{equation}
with value zero on the exceptional null sets, has expectation
one under the law that draws $B$ and projects independent $p$-distributed
signals. Conditional on the replica tuple and its original matrix $A$,
$\widetilde A$ has law $\gamma_{k,d}^{F}$; this law depends only on the
tuple. If $K'$ is the projected total correlation of the replicas, then
\begin{equation}\label{eq:gaussian-correlation-loss}
 K-K'\le Ckm+
 m\E\!\left[\log p_{\widetilde A}(\widetilde A S_1)-\log p_A(Y)\right].
\end{equation}
Every logarithm in this display is absolutely integrable.
\end{lemma}

\begin{proof}
For a fixed full-rank tuple write $D=ET$, where the columns of $E$ form an
orthonormal basis of $F$ and $T$ is invertible. Then $G=BE$ is standard
Gaussian and
\[
 \det((BD)^{\mathsf T}(BD))
 =J^2\det(G^{\mathsf T}G).
\]
The action of $B$ on $F^\perp$ is independent of $G$. Conditional on $G$
and the tuple, $C_0$ is fixed with orthonormal rows and $C_0BE=0$.
Thus $\widetilde A=C_0B$ has law $\gamma_{k,d}^{F}$, independently of
$G$ given the tuple. Since $B$ is independent of the replica experiment,
this remains its law after conditioning also on the original $A$.
The test is observable from $(B,U_1,\ldots,U_t)$:
$C_0$, $C_0B$, and its evaluated label $C_0U_1$ are all determined there.
The common density versions of Theorem~\ref{thm:equal-label-measure} make
the last factor measurable, finite, and positive almost surely under the
independent-tuple reference. The replica tuple is absolutely continuous
with respect to that reference's tuple marginal, so the same assertion
holds for replicas.

Under the independent-signal reference law, integrate the factor
$\det(G^{\mathsf T}G)^{-k/2}$ first. It contributes $Z_{k,m}$.
The remaining expectation is
\[
 \int cJ^{-k}p_{\widetilde A}(\widetilde A S_1)^{-m}
       p^{\otimes t}(d\mathbf s)\gamma_{k,d}^{F}(d\widetilde A)=1
\]
by Corollary~\ref{cor:replica-likelihood}. This proves the normalization
of \eqref{eq:gaussian-projected-test}.

Here are the needed Gaussian determinant bounds. Gram--Schmidt expresses
$\det(G^{\mathsf T}G)$ as a product of independent chi-squares with degrees
$\nu=\ell,\ell-1,\ldots,\ell-m+1$; this standard Gaussian reduction is also
described by Edelman and Rao~\cite[Section 5]{EdelmanRao2005}.
The smallest degree is $3k+1>k$. Since $k$ is even, direct integration of
the chi-square density gives
\[
 \E(\chi_\nu^2)^{-k/2}
 =\prod_{j=1}^{k/2}(\nu-2j)^{-1}\le(2k)^{-k/2}.
\]
In particular $Z_{k,m}<\infty$. The same densities give absolute
logarithmic integrability, and Jensen gives
$\E\log\det(G^{\mathsf T}G)\le m\log\ell$. Consequently
\begin{equation}\label{eq:gaussian-wishart-remainder}
 \log Z_{k,m}+\frac{k}{2}\E\log\det(G^{\mathsf T}G)\le Ckm.
\end{equation}

Before taking logarithms of the test, we control the positive logarithm
of its projection density. For a fixed tuple put
\[
 \delta(s')=\norm{P_{F^\perp}(s'-S_1)}.
\]
The ball-average density version in
Theorem~\ref{thm:equal-label-measure} holds at
$\widetilde A S_1$ under the independent-tuple reference and hence under
the replica law. The coordinates of $\widetilde A(s'-S_1)$ are
independent centered Gaussians with variance $\delta(s')^2$.
A ball has probability at most its volume times the maximum Gaussian
density. Apply this bound to the ball averages and then use Fatou's lemma:
\begin{equation}\label{eq:gaussian-negative-moment}
 \E_{\widetilde A}p_{\widetilde A}(\widetilde A S_1)
 \le (2\pi)^{-k/2}\int\delta(s')^{-k}\,p(ds').
\end{equation}
Here and below such conditional estimates hold for almost every tuple.
The affine space $S_1+F$ has dimension $m$, so
\eqref{eq:affine-tube} gives
\[
 p\{\delta\le r\}\le LC^d r^{d-1-m},\qquad 0<r\le1.
\]
Since $d-1-m>k$, tail integration bounds the inverse moment uniformly
over tuples for fixed $d,L$. Thus
$\log^+p_{\widetilde A}(\widetilde A S_1)$ is integrable.

The projected law has divergence $K'\le K<\infty$ from the
independent-signal reference. Normalization of the test and
\eqref{eq:entropy-positive-log} imply that $(\log\Lambda)^+$ is
integrable. Its defining identity is
\begin{equation}\label{eq:gaussian-test-log}
 \begin{aligned}
 \log\Lambda
 &=\log c-\log Z_{k,m}-k\log J
   -\frac{k}{2}\log\det(G^{\mathsf T}G)\\
 &\qquad-m\log p_{\widetilde A}(\widetilde A S_1).
 \end{aligned}
\end{equation}
The negative part of this expression is bounded by a constant and the
positive logarithms of $J$, the Gaussian determinant, and the last
density. Those are integrable by Lemma~\ref{lem:replica-logs}, the
chi-square calculation, and \eqref{eq:gaussian-negative-moment}.
Thus $\log\Lambda$ is absolutely integrable. Solving
\eqref{eq:gaussian-test-log} for the last logarithm proves its absolute
integrability as well.

We may now apply the entropy test \eqref{eq:entropy-test} to obtain
$K'\ge\E\log\Lambda$. Subtract this bound from
\eqref{eq:gaussian-correlation-upper}. Equation
\eqref{eq:gaussian-test-log} displays the exact cancellation of
$-k\E\log J$, and \eqref{eq:gaussian-wishart-remainder} bounds the
remaining determinant terms by $Ckm$. This gives
\eqref{eq:gaussian-correlation-loss}.
\end{proof}

\subsection{Comparing the remaining projection densities}

Both the actual $A$ and the observable $\widetilde A$ annihilate $F$.
The preceding lemma has removed the affine volume from the information
bound. It remains to compare the two densities in
\eqref{eq:gaussian-correlation-loss}. We first bound the conditional mean
of the new density by a sum of averages of the actual density. The radii
in that sum will depend only on $A$, allowing us to average the actual
label with its original density $p_A$.

\begin{lemma}[Dyadic density comparison]\label{lem:gaussian-density}
In the replica experiment with the independent side matrix and
$\widetilde A=C_0B$ defined above,
\begin{equation}\label{eq:gaussian-density-comparison}
 \E\!\left[\log p_{\widetilde A}(\widetilde A S_1)-\log p_A(Y)\right]
 \le Ck+C\log(2+\log L).
\end{equation}
\end{lemma}

\begin{proof}
Fix the tuple and its actual $(A,Y)$, and again write
$\delta(s')=\norm{P_{F^\perp}(s'-S_1)}$. Conditional on these data,
$\widetilde A$ has law $\gamma_{k,d}^{F}$. We refine the inverse-moment
bound \eqref{eq:gaussian-negative-moment} by using $A$.
Choose the deterministic integer
\begin{equation}\label{eq:gaussian-cutoff}
 J_0=1+\left\lceil
       \frac{\log_2 L+C_1d}{d-1-m-k}\right\rceil
\end{equation}
with a sufficiently large absolute $C_1$. The denominator is bounded
below by a positive absolute multiple of $d$. The contribution of
$\delta\le2^{-J_0}$ to the inverse moment, including the Gaussian constant,
is at most
\[
 (2\pi)^{-k/2}
 \sum_{j\ge J_0}2^{(j+1)k}LC^d2^{-j(d-1-m)}\le1.
\]
For the remaining distances use the bins
$\delta\in(r/2,r]$, where $r=2,1,\ldots,2^{-J_0+1}$.
These cover the remaining values because $\delta(s')\le\norm{s'-S_1}\le2$.
Since $A$ annihilates $F$,
\[
 \norm{As'-Y}\le\norm A_{\op}\,\delta(s').
\]
Let $v_k$ be the unit $k$-ball volume, and let
$\overline p_A^{\,a}(y)$ be the average of $p_A$ over the radius-$a$ ball
centered at $y$. The contribution of one bin is at most
\[
 (2/r)^k p\{\delta\le r\}
 \le 2^k v_k\norm A_{\op}^{k}
       \overline p_A^{\,\norm A_{\op}r}(Y).
\]
Integrating $e^{-\norm z^2/2}$ on the radius-$\sqrt{k}$ ball gives
$v_k\le(2\pi e/k)^{k/2}$. Hence
\begin{equation}\label{eq:gaussian-density-sum}
 \E_{\widetilde A}p_{\widetilde A}(\widetilde A S_1)
 \le1+\left(\frac{C\norm A_{\op}}{\sqrt{k}}\right)^k
       \sum_r\overline p_A^{\,\norm A_{\op}r}(Y).
\end{equation}
The right side depends on the tuple only through $(A,Y)$.

Divide by the finite positive density $p_A(Y)$ and apply Jensen first to
$\widetilde A$; all logarithms are integrable by
Lemma~\ref{lem:gaussian-test}. Only after obtaining
\eqref{eq:gaussian-density-sum} do we average over $Y$ given $A$.
Its law is then the original $p_A(y)\,dy$, and for every radius fixed by $A$,
\[
 \E_{Y\mid A}\frac{\overline p_A^{\,a}(Y)}{p_A(Y)}\le1,\qquad
 \E_{Y\mid A}\frac1{p_A(Y)}
 \le v_k\norm A_{\op}^{k}.
\]
The first inequality integrates an averaged probability density over
$\{p_A>0\}$, and the second is the volume of that set, contained in the
radius-$\norm A_{\op}$ ball. A second application of Jensen thus bounds
the expected log ratio by
\begin{equation}\label{eq:gaussian-density-jensen}
 k\E\log\!\left(C\left(1+\frac{\norm A_{\op}}{\sqrt{k}}\right)\right)
       +\log(J_0+2).
\end{equation}
This order of averaging is essential: before the tuple has disappeared,
its conditional label need not have density $p_A$ under further conditioning.

For completeness, choose $1/4$-nets of the unit spheres in $\R^d$ and
$\R^k$, with at most $9^d$ and $9^k$ points. The operator norm is at most
twice the maximum absolute Gaussian bilinear form on these nets. The
expected maximum of $N$ absolute standard normals is at most
$\sqrt{2\log(2N)}$, by the exponential-moment bound; their independence is
not required. Thus $\E\norm A_{\op}\le C\sqrt d$. Since $k$ is a fixed
positive fraction of $d$, Jensen bounds the first term of
\eqref{eq:gaussian-density-jensen} by $Ck$. The cutoff
\eqref{eq:gaussian-cutoff} gives
$\log(J_0+2)\le C\log(2+\log L)$, proving the lemma.
\end{proof}

\begin{proof}[Proof of Theorem~\ref{thm:gaussian-block}]
Generate $W$ conditionally on $(A,Y)$ independently of the replicas.
Each pair $(S_i,W)$ has the original pair law, and the side matrix $B$ is
independent of the whole construction. Lemma~\ref{lem:replica-information},
followed by Lemmas~\ref{lem:gaussian-test} and
\ref{lem:gaussian-density}, gives
\[
 tI(S;W\mid B,BS)
 \le H(W)+Ckm+m\bigl(Ck+C\log(2+\log L)\bigr).
\]
Here the representative signal is coupled independently with $B$, exactly
as in the original block experiment. Since $t=m+1$ and $k$ is proportional
to $d$, division by $t$ proves \eqref{eq:gaussian-block}.
\end{proof}
